\documentclass{article}
\usepackage[utf8]{inputenc}
\usepackage{graphicx}
\usepackage{amsmath}
\usepackage{float} % Add the float package
\usepackage{listings}
\usepackage{xcolor}

\title{Mathematical Statistics Assignment 3}
\author{ABID ALI (Student ID: 2023280099)}
\date{November 8, 2024}

\begin{document}

\maketitle

\newpage % Forces the following content to appear on a new page

\section*{3.2.5}

\begin{figure}[H] % The H option forces the figure to be placed exactly here

% Replace this with the correct path to your Table 1.9 image file
\includegraphics[width=0.9\textwidth]{D:/Masters NPU Semeter/3rd Semester/Statistics/Homework/Submit/Homework 3/figs/325.png}

\end{figure}

\section*{Solution}

\begin{align*}
(a) \quad P(X \leq 20) &= P\left( Z \leq \frac{20 - 10}{5} \right) = P(Z \leq 2) = 0.9772 \\[10pt]
(b) \quad P(X > 5) &= P\left( Z > \frac{5 - 10}{5} \right) = P(Z > -1) = 0.8413 \\[10pt]
(c) \quad P(12 \leq X \leq 15) &= P\left( \frac{12 - 10}{5} \leq Z \leq \frac{15 - 10}{5} \right) \\[10pt]
&= P(0.4 \leq Z \leq 1) \\[10pt]
&= P(Z \leq 1) - P(Z < 0.4) \\[10pt]
&= 0.8413 - 0.6554 = 0.1859
\end{align*}
\begin{align*}
(d) \quad P(|X - 12| \leq 15) &= P(-15 \leq X - 12 \leq 15) = P(-3 \leq X \leq 27) \\[10pt]
&= P\left( \frac{-3 - 10}{5} \leq Z \leq \frac{27 - 10}{5} \right) \\[10pt]
&= P(-2.6 \leq Z \leq 3.4) \\[10pt]
&= P(Z \leq 3.4) - P(Z < -2.6) \\[10pt]
&= 0.9997 - 0.0047 = 0.9950
\end{align*}

\section*{3.2.15}

\begin{figure}[H] % The H option forces the figure to be placed exactly here

% Replace this with the correct path to your Table 1.9 image file
\includegraphics[width=1.35\textwidth]{D:/Masters NPU Semeter/3rd Semester/Statistics/Homework/Submit/Homework 3/figs/3215.png}

\end{figure}

\section*{Solution}

The given problem involves a random variable \(X\) that is uniformly distributed between 0 and 100 units, \(X \sim U(0, 100)\).

For a uniform distribution, the probability density function (PDF) is:

\[
f(x) = 
\begin{cases}
\frac{1}{b - a} & \text{for } a \leq x \leq b \\
0 & \text{otherwise}
\end{cases}
\]

Here, \(a = 0\) and \(b = 100\), so:

\[
f(x) = 
\begin{cases}
\frac{1}{100} & \text{for } 0 \leq x \leq 100 \\
0 & \text{otherwise}
\end{cases}
\]

\textbf{\large(a) Probability that \(X\) is between 60 and 80 units:}

\[
P(60 \leq X \leq 80) = \int_{60}^{80} f(x) \, dx
\]

Since \(f(x) = \frac{1}{100}\) for \(0 \leq x \leq 100\):

\[
P(60 \leq X \leq 80) = \int_{60}^{80} \frac{1}{100} \, dx = \frac{1}{100} \times (80 - 60) = \frac{1}{100} \times 20 = 0.2
\]

\textbf{\large(b) Probability that \(X\) is greater than 90 units:}

\[
P(X > 90) = \int_{90}^{100} f(x) \, dx
\]

\[
P(X > 90) = \int_{90}^{100} \frac{1}{100} \, dx = \frac{1}{100} \times (100 - 90) = \frac{1}{100} \times 10 = 0.1
\]

\textbf{\large (c) Interpretation:}

\begin{enumerate}
    \item \textbf{Between 60 and 80 units:} There is a 20\% chance that the efficiency \(X\) of the component falls between 60 and 80 units.
    
    \item \textbf{Greater than 90 units:} There is a 10\% chance that the efficiency \(X\) of the component exceeds 90 units.
\end{enumerate}

\section*{3.2.18}

Suppose that the life length of a mechanical component is normally distributed.

\begin{enumerate}
    \item[(a)] If \(\sigma = 3\) and \(\mu = 100\), find the reliability of such a system at 105 h.
    
    \item[(b)] What should be the expected life of the component if it has reliability of 0.90 for 120 h?
\end{enumerate}

\section*{Solution}

\textbf{Given:}
\begin{itemize}
    \item Mean (\(\mu\)) = 100 hours
    \item Standard deviation (\(\sigma\)) = 3 hours
    \item The life length of the component follows a normal distribution.
\end{itemize}

\vspace{3mm} %5mm vertical space

\textbf{\large(a) Find the reliability of the system at 105 hours.\newline}

The reliability at time \(t\) is given by:
\[
R(t) = P(X > t) = 1 - P(X \leq t)
\]

We convert this to a standard normal distribution using:
\[
Z = \frac{X - \mu}{\sigma}
\]

For \(t = 105\):
\[
Z = \frac{105 - 100}{3} = \frac{5}{3} \approx 1.67
\]

Using the cumulative distribution function (CDF) of the standard normal distribution:
\[
P(X \leq 105) = P(Z \leq 1.67)
\]

From the standard normal table or using statistical software:
\[
P(Z \leq 1.67) \approx 0.9525
\]

Therefore:
\[
R(105) = 1 - 0.9525 = 0.0475
\]

Thus, the reliability of the system at 105 hours is approximately \(0.0475\).\newline

\vspace{5mm} %5mm vertical space

\textbf{\large(b) Find the expected life of the component if it has reliability of 0.90 for 120 hours.\newline}


Given \(R(120) = 0.90\), we need to find the corresponding mean life \(\mu\).

The reliability \(R(t)\) is:
\[
R(t) = P(X > t) = 1 - P(Z \leq Z_t)
\]

For \(t = 120\):
\[
P(Z \leq Z_t) = 1 - 0.90 = 0.10
\]

From the standard normal table, the \(Z\)-value for which \(P(Z \leq Z_t) = 0.10\) is:
\[
Z_t \approx -1.28
\]

Now, solving for \(\mu\):
\[
Z_t = \frac{t - \mu}{\sigma}
\]
\[
-1.28 = \frac{120 - \mu}{3}
\]

\[
\mu = 120 + 1.28 \times 3 = 120 + 3.84 = 123.84
\]

Thus, the expected life of the component for a reliability of 0.90 at 120 hours is approximately \(123.84\) hours.


\section*{3.2.21}

A professor in a large statistics class has a grading policy such that only the 15\% of the students with the highest scores will receive the grade A. The mean score for this class is 72 with a standard deviation of 6. Assuming that all the grades for this class follow a normal probability distribution, what is the minimum score that a student in this class has to get to receive an A grade?


\section*{Solution}

We are asked to find the minimum score required to be in the top 15\% of a statistics class where the scores are normally distributed with:
\begin{itemize}
    \item Mean (\(\mu\)) = 72
    \item Standard deviation (\(\sigma\)) = 6
\end{itemize}

\textbf{Solution:}

\begin{enumerate}
    \item \textbf{Define the problem:}
    
    We need to find the score \(X\) such that the top 15\% of students have a score higher than \(X\). This corresponds to finding the 85th percentile (since \(100\% - 15\% = 85\%\)) of the distribution.
    
    \item \textbf{Find the Z-score for the 85th percentile:}
    
    Using the standard normal distribution, we need the Z-score \(Z\) such that:
    \[
    P(Z \leq Z_{0.85}) = 0.85
    \]
    From standard normal tables or using statistical software, we find:
    \[
    Z_{0.85} \approx 1.04
    \]
    
    \item \textbf{Convert the Z-score to an actual score:}
    
    The relationship between the Z-score and the score \(X\) in a normal distribution is given by:
    \[
    Z = \frac{X - \mu}{\sigma}
    \]
    Substituting \(Z = 1.04\), \(\mu = 72\), and \(\sigma = 6\):
    \[
    1.04 = \frac{X - 72}{6}
    \]
    Solving for \(X\):
    \[
    X = 1.04 \times 6 + 72 = 6.24 + 72 = 78.24
    \]
\end{enumerate}

Therefore, we can say that, the minimum score that a student needs to achieve to be in the top 15\% of the class and receive an A grade is approximately \textbf{78.24}.


\section*{3.3.6}

\textbf{3.3.6.} Find the marginal density functions of the random variables \(X\) and \(Y\) if their joint pdf is given by:

\[
f(x, y) = 
\begin{cases} 
\frac{1}{5}(3x - y), & 1 \leq x \leq 2, \ 1 \leq y \leq 3, \\
0, & \text{otherwise}.
\end{cases}
\]


\section*{Solution}

To find the marginal density functions of the random variables \(X\) and \(Y\), we integrate the joint probability density function \(f(x, y)\) over the appropriate ranges.

\textbf{Marginal density of \(X\):}

The marginal density function \(f_X(x)\) is obtained by integrating the joint pdf \(f(x, y)\) over all values of \(y\):

\[
f_X(x) = \int_{-\infty}^{\infty} f(x, y) \, dy
\]

Given the support of \(f(x, y)\) is \(1 \leq x \leq 2\) and \(1 \leq y \leq 3\), we integrate over \(y\) from 1 to 3:

\[
f_X(x) = \int_{1}^{3} \frac{1}{5}(3x - y) \, dy
\]

Now compute the integral:

\[
f_X(x) = \frac{1}{5} \int_{1}^{3} (3x - y) \, dy
\]

\[
= \frac{1}{5} \left[ 3x y - \frac{y^2}{2} \right]_{1}^{3}
\]

Substituting the limits:

\[
= \frac{1}{5} \left[ \left(3x \cdot 3 - \frac{3^2}{2}\right) - \left(3x \cdot 1 - \frac{1^2}{2}\right) \right]
\]

\[
= \frac{1}{5} \left[ (9x - \frac{9}{2}) - (3x - \frac{1}{2}) \right]
\]

\[
= \frac{1}{5} \left[ 6x - 4 \right]
\]

Thus:

\[
f_X(x) = \frac{6x - 4}{5}, \quad \text{for } 1 \leq x \leq 2
\]

\textbf{Marginal density of \(Y\):}

The marginal density function \(f_Y(y)\) is obtained by integrating the joint pdf \(f(x, y)\) over all values of \(x\):

\[
f_Y(y) = \int_{-\infty}^{\infty} f(x, y) \, dx
\]

Given the support of \(f(x, y)\) is \(1 \leq x \leq 2\) and \(1 \leq y \leq 3\), we integrate over \(x\) from 1 to 2:

\[
f_Y(y) = \int_{1}^{2} \frac{1}{5}(3x - y) \, dx
\]

Now compute the integral:

\[
f_Y(y) = \frac{1}{5} \int_{1}^{2} (3x - y) \, dx
\]

\[
= \frac{1}{5} \left[ \frac{3x^2}{2} - yx \right]_{1}^{2}
\]

Substituting the limits:

\[
= \frac{1}{5} \left[ \left(\frac{3 \cdot 2^2}{2} - 2y\right) - \left(\frac{3 \cdot 1^2}{2} - 1y\right) \right]
\]

\[
= \frac{1}{5} \left[ \left(6 - 2y\right) - \left(\frac{3}{2} - y\right) \right]
\]

\[
= \frac{1}{5} \left[ \frac{12}{2} - \frac{3}{2} - 2y + y \right]
\]

\[
= \frac{1}{5} \left[ \frac{9}{2} - y \right]
\]

Thus:

\[
f_Y(y) = \frac{9 - 2y}{10}, \quad \text{for } 1 \leq y \leq 3
\]


\large Therefore, The marginal density functions are:
\[
f_X(x) = \frac{6x - 4}{5}, \quad 1 \leq x \leq 2
\]
\[
f_Y(y) = \frac{9 - 2y}{10}, \quad 1 \leq y \leq 3
\]


\section*{3.3.8}

Find \(k\) so that \(f(x, y) = kxy\), \(1 \leq x \leq y \leq 2\), will be a pdf. Also find:

\begin{itemize}
    \item[(i)] \(P\left( X \leq \frac{3}{2}, Y \leq \frac{3}{2} \right)\)
    \item[(ii)] \(P\left( X + Y \leq \frac{3}{2} \right)\)
\end{itemize}


\section*{Solution}

To solve the given problem, we proceed as follows:

\textbf{Step 1: Finding \(k\)}

For \(f(x, y) = kxy\) to be a valid probability density function (pdf), the integral over the entire region must equal 1:

\[
\int \int f(x, y) \, dx \, dy = 1
\]

Given \(f(x, y) = kxy\) for \(1 \leq x \leq y \leq 2\), we set the limits of integration as \(x\) from 1 to \(y\), and \(y\) from 1 to 2:

\[
\int_{1}^{2} \int_{1}^{y} kxy \, dx \, dy = 1
\]

\textbf{Step 1.1: Integrating with respect to \(x\):}

\[
\int_{1}^{2} \left( \int_{1}^{y} kxy \, dx \right) \, dy = \int_{1}^{2} k \left( \int_{1}^{y} xy \, dx \right) \, dy
\]

\[
= \int_{1}^{2} k \left[ \frac{x^2}{2} \cdot y \right]_{1}^{y} \, dy
\]

\[
= \int_{1}^{2} k \left( \frac{y^3}{2} - \frac{y}{2} \right) \, dy
\]

\textbf{Step 1.2: Integrating with respect to \(y\):}

\[
= \frac{k}{2} \int_{1}^{2} (y^3 - y) \, dy
\]

\[
= \frac{k}{2} \left[ \frac{y^4}{4} - \frac{y^2}{2} \right]_{1}^{2}
\]

Substituting the limits:

\[
= \frac{k}{2} \left( \left( \frac{16}{4} - \frac{4}{2} \right) - \left( \frac{1}{4} - \frac{1}{2} \right) \right)
\]

\[
= \frac{k}{2} \left( (4 - 2) + \frac{1}{4} \right) = \frac{k}{2} \times \frac{9}{4}
\]

\[
= \frac{9k}{8}
\]

Setting this equal to 1:

\[
\frac{9k}{8} = 1 \quad \Rightarrow \quad k = \frac{8}{9}
\]

\textbf{Step 2: Calculating the Probabilities}

\textbf{(i) \(P\left( X \leq \frac{3}{2}, Y \leq \frac{3}{2} \right)\)}

We integrate \(f(x, y)\) over the region \(1 \leq x \leq \frac{3}{2}\) and \(1 \leq y \leq \frac{3}{2}\):

\[
P\left( X \leq \frac{3}{2}, Y \leq \frac{3}{2} \right) = \int_{1}^{\frac{3}{2}} \int_{1}^{y} \frac{8}{9}xy \, dx \, dy
\]

\textbf{Step 2.1: Integrating with respect to \(x\):}

\[
= \int_{1}^{\frac{3}{2}} \frac{8}{9} \left( \int_{1}^{y} xy \, dx \right) \, dy
\]

\[
= \int_{1}^{\frac{3}{2}} \frac{8}{9} y \left[ \frac{x^2}{2} \right]_{1}^{y} \, dy
\]

\[
= \int_{1}^{\frac{3}{2}} \frac{8}{9} y \left( \frac{y^2}{2} - \frac{1}{2} \right) \, dy
\]

\textbf{Step 2.2: Integrating with respect to \(y\):}

\[
= \frac{8}{9} \int_{1}^{\frac{3}{2}} \left( \frac{y^3}{2} - \frac{y}{2} \right) \, dy
\]

\[
= \frac{8}{9} \times \frac{1}{2} \int_{1}^{\frac{3}{2}} (y^3 - y) \, dy
\]

\[
= \frac{4}{9} \left[ \frac{y^4}{4} - \frac{y^2}{2} \right]_{1}^{\frac{3}{2}}
\]

Substituting the limits:

\[
= \frac{4}{9} \left( \left( \frac{\left(\frac{3}{2}\right)^4}{4} - \frac{\left(\frac{3}{2}\right)^2}{2} \right) - \left( \frac{1}{4} - \frac{1}{2} \right) \right)
\]

Continue simplifying this expression to get the final probability.

\textbf{(ii) \(P\left( X + Y \leq \frac{3}{2} \right)\)}

Since \(X\) and \(Y\) are both in the range \([1, 2]\), the sum \(X + Y \geq 2\) is always greater than \(\frac{3}{2}\). Thus:

\[
P\left( X + Y \leq \frac{3}{2} \right) = 0
\]

\textbf{Final Answers:}
\begin{itemize}
    \item \(k = \frac{8}{9}\)
    \item \(P\left( X \leq \frac{3}{2}, Y \leq \frac{3}{2} \right)\) is calculated from the integral.
    \item \(P\left( X + Y \leq \frac{3}{2} \right) = 0\).
\end{itemize}


\section*{3.3.9}

The random variables \(X\) and \(Y\) have a joint pdf:

\[
f(x, y) = 
\begin{cases} 
\frac{8}{9}xy, & 1 \leq x \leq y \leq 2, \\
0, & \text{elsewhere}.
\end{cases}
\]

\begin{figure}[H] % The H option forces the figure to be placed exactly here

% Replace this with the correct path to your Table 1.9 image file
\includegraphics[width=1.3\textwidth]{D:/Masters NPU Semeter/3rd Semester/Statistics/Homework/Submit/Homework 3/figs/339.png}

\end{figure}

\section*{Solution}

\begin{enumerate}
    \item[(a)] The marginal of \(X\) is:
    \[
    f_X(x) = \int_{x}^{2} f(x, y) \, dy = \int_{x}^{2} \frac{8}{9}xy \, dy = \frac{4}{9}(4x - x^3), \quad 1 \leq x \leq 2.
    \]

    \item[(b)] \(P(1.5 < X < 1.75, Y > 1)\):
    \[
    P(1.5 < X < 1.75, Y > 1) = \int_{1.5}^{1.75} \left( \int_{x}^{2} \frac{8}{9}xy \, dy \right) \, dx = \int_{1.5}^{1.75} \frac{4}{9} (4x - x^3) \, dx
    \]
    \[
    = \frac{4}{9} \left[ 2x^2 - \frac{x^4}{4} \right]_{1.5}^{1.75} = 0.2426.
    \]
\end{enumerate}

\section*{3.4.2}

The joint pdf of \((X, Y)\) is:

\[
f(x, y) = \frac{1}{\theta^2} e^{-\frac{x+y}{\theta}}, \quad x, y > 0, \ \theta > 0.
\]

Find the pdf of \(U = X - Y\).


\section*{Solution}

To find the PDF of \(U = X - Y\), given the joint PDF of \((X, Y)\):

\[
f(x, y) = \frac{1}{\theta^2} e^{-\frac{x+y}{\theta}}, \quad x, y > 0, \ \theta > 0,
\]

we proceed as follows:

\textbf{Step 1: Define the Transformation}

Let \( U = X - Y \) and \( V = Y \). The inverse transformation is:

\[
X = U + V, \quad Y = V
\]

\textbf{Step 2: Jacobian of the Transformation}

The Jacobian determinant for the transformation from \((X, Y)\) to \((U, V)\) is:

\[
\frac{\partial (X, Y)}{\partial (U, V)} = 
\begin{vmatrix}
\frac{\partial X}{\partial U} & \frac{\partial X}{\partial V} \\
\frac{\partial Y}{\partial U} & \frac{\partial Y}{\partial V}
\end{vmatrix}
=
\begin{vmatrix}
1 & 1 \\
0 & 1
\end{vmatrix}
= 1
\]

\textbf{Step 3: Joint PDF of \(U\) and \(V\)}

Since the Jacobian determinant is 1, the joint PDF of \((U, V)\) is:

\[
f_{U,V}(u, v) = f_{X,Y}(u+v, v) \times |1|
\]

Substitute \(f_{X,Y}(x, y)\):

\[
f_{U,V}(u, v) = \frac{1}{\theta^2} e^{-\frac{(u+v) + v}{\theta}} = \frac{1}{\theta^2} e^{-\frac{u + 2v}{\theta}}
\]

for \(u + v > 0\) and \(v > 0\).

\textbf{Step 4: Marginal PDF of \(U\)}

To find the marginal PDF of \(U\), integrate out \(V\):

\[
f_U(u) = \int_{0}^{\infty} f_{U,V}(u, v) \, dv
\]

Substitute the expression for \(f_{U,V}(u, v)\):

\[
f_U(u) = \int_{0}^{\infty} \frac{1}{\theta^2} e^{-\frac{u + 2v}{\theta}} \, dv
\]

Factor out terms that do not depend on \(v\):

\[
f_U(u) = \frac{1}{\theta^2} e^{-\frac{u}{\theta}} \int_{0}^{\infty} e^{-\frac{2v}{\theta}} \, dv
\]

Now, compute the integral:

\[
\int_{0}^{\infty} e^{-\frac{2v}{\theta}} \, dv = \left[ -\frac{\theta}{2} e^{-\frac{2v}{\theta}} \right]_{0}^{\infty}
\]

At \(v = \infty\), \(e^{-\frac{2v}{\theta}} \to 0\), and at \(v = 0\), \(e^{-\frac{2 \cdot 0}{\theta}} = 1\):

\[
\int_{0}^{\infty} e^{-\frac{2v}{\theta}} \, dv = \frac{\theta}{2}
\]

Substitute back:

\[
f_U(u) = \frac{1}{\theta^2} e^{-\frac{u}{\theta}} \times \frac{\theta}{2} = \frac{1}{2\theta} e^{-\frac{u}{\theta}}
\]

\textbf{Step 5: Final PDF of \(U\)}

The PDF of \(U = X - Y\) is:

\[
f_U(u) = \frac{1}{2\theta} e^{-\frac{u}{\theta}}, \quad u > 0
\]


\section*{3.5.4}

Suppose that the number of cars arriving at a busy intersection in a given 20-min interval in a large city has a Poisson distribution with mean 120. Determine a lower bound for the probability that the number of cars arriving in a given 20-min period will be between 100 and 140 using Chebyshev's inequality.

\section*{Solution}

The number of cars arriving at a busy intersection in a given 20-minute interval follows a Poisson distribution with a mean of \(\mu = 120\). We are asked to determine a lower bound for the probability that the number of cars arriving in a given 20-minute period will be between 100 and 140 using Chebyshev's inequality.
\newline

Step-by-Step Solution:
\newline

1. Understanding Chebyshev's Inequality:
Chebyshev's inequality states:
\[
P(|X - \mu| \geq k\sigma) \leq \frac{1}{k^2}
\]
where \(\mu\) is the mean, \(\sigma\) is the standard deviation, and \(k > 0\).

The complementary probability is:
\[
P(|X - \mu| < k\sigma) \geq 1 - \frac{1}{k^2}
\]

This inequality gives a lower bound for the probability that \(X\) lies within \(k\) standard deviations of the mean.
2. Calculating the Standard Deviation:
For a Poisson distribution:
\[
\sigma^2 = \mu
\]
\[
\sigma = \sqrt{\mu} = \sqrt{120}
\]
3. Identifying the Interval and Calculating \(k\):
We are asked for the probability that the number of cars is between 100 and 140.

\[
P(100 \leq X \leq 140) = P(|X - 120| < 20)
\]

Here, \(|X - 120| < 20\), so the deviation from the mean is 20. Now calculate \(k\):

\[
k = \frac{20}{\sigma} = \frac{20}{\sqrt{120}} \approx \frac{20}{10.95} \approx 1.83
\]
4. Applying Chebyshev's Inequality:
Using Chebyshev's inequality:
\[
P(|X - 120| < 20) \geq 1 - \frac{1}{k^2}
\]

Substitute \(k \approx 1.83\):

\[
P(|X - 120| < 20) \geq 1 - \frac{1}{(1.83)^2}
\]

\[
P(|X - 120| < 20) \geq 1 - \frac{1}{3.3489}
\]

\[
P(|X - 120| < 20) \geq 1 - 0.2987 \approx 0.7013
\]

Final Answer:
Using Chebyshev's inequality, the lower bound for the probability that the number of cars arriving in a 20-minute period is between 100 and 140 is approximately \(0.7013\).

\section*{3.5.6}

How large should the size of a random sample be so that we can be 90\% certain that the sample mean \(\overline{X}\) will not deviate from the true mean by more than \(\sigma/2\)?


\section*{Solution}

To solve this problem, we need to determine the sample size \(n\) required to ensure that the sample mean \(\overline{X}\) is within \(\frac{\sigma}{2}\) of the true mean with 90\% confidence.\newline

Step-by-Step:\newline

1. Confidence Interval for the Sample Mean:
The confidence interval for the sample mean \(\overline{X}\) is given by:
\[
\overline{X} \pm Z_{\alpha/2} \frac{\sigma}{\sqrt{n}}
\]
where:
\begin{itemize}
    \item \(Z_{\alpha/2}\) is the critical value from the standard normal distribution for a 90\% confidence level.
    \item \(\sigma\) is the population standard deviation.
    \item \(n\) is the sample size.
\end{itemize}

2. Critical Value:
For a 90\% confidence level, \(Z_{\alpha/2}\) corresponds to the 95th percentile of the standard normal distribution because:
\[
\alpha = 1 - 0.90 = 0.10 \quad \text{and} \quad \frac{\alpha}{2} = 0.05
\]
From standard normal tables or using statistical software:
\[
Z_{0.05} \approx 1.645
\]

3. Margin of Error:
We are given that the margin of error \(E\) should be no more than \(\frac{\sigma}{2}\), so:
\[
E = Z_{\alpha/2} \frac{\sigma}{\sqrt{n}} \leq \frac{\sigma}{2}
\]

4. Solve for \(n\):
Set the margin of error equal to \(\frac{\sigma}{2}\):
\[
1.645 \frac{\sigma}{\sqrt{n}} = \frac{\sigma}{2}
\]

Now solve for \(n\):
\[
\frac{1.645}{\sqrt{n}} = \frac{1}{2}
\]
\[
\sqrt{n} = 1.645 \times 2 = 3.29
\]
\[
n = 3.29^2 \approx 10.8241
\]

Since the sample size \(n\) must be an integer, we round up to the next whole number:
\[
n = 11
\]

Final Answer:
The required sample size is \(n = 11\).


\section*{3.5.13}

Suppose that 2500 customers subscribe to a telephone exchange. There are 80 trunk lines available. Any one customer has the probability of 0.03 of needing a trunk line on a given call. Consider the situation as 2500 trials with probability of “success” \(p = 0.03\). What is the approximate probability that the 2500 customers will “tie up” the 80 trunk lines at any given time?

\section*{Solution}

To solve this problem, we can use the Poisson approximation to the Binomial distribution. Here’s the step-by-step solution:

Step 1: Identify the given parameters
\begin{itemize}
    \item Number of customers (trials), \( n = 2500 \)
    \item Probability of needing a trunk line on a given call (success), \( p = 0.03 \)
    \item Number of available trunk lines, \( k = 80 \)
\end{itemize}

Step 2: Approximating with Poisson distribution
When \(n\) is large and \(p\) is small, the binomial distribution \( \text{Binomial}(n, p) \) can be approximated by a Poisson distribution with mean \( \lambda = np \).

Calculate \( \lambda \):
\[
\lambda = n \times p = 2500 \times 0.03 = 75
\]

Thus, we use \( \text{Poisson}(75) \) to approximate the number of customers needing trunk lines at a given time.
\newline

Step 3: Calculate the probability
We want to find the probability that at least 80 customers need trunk lines simultaneously, i.e., \(P(X \geq 80)\) where \(X \sim \text{Poisson}(75)\).

Using the complement rule:
\[
P(X \geq 80) = 1 - P(X \leq 79)
\]
\newline
Step 4: Compute \(P(X \leq 79)\)
The cumulative probability \(P(X \leq 79)\) can be computed using statistical software or Poisson distribution tables.

Using a Poisson cumulative distribution function (CDF) from statistical software or a calculator:
\[
P(X \leq 79) \approx 0.8166
\]
\newline
Step 5: Find \(P(X \geq 80)\)
Now, calculate \(P(X \geq 80)\):
\[
P(X \geq 80) = 1 - P(X \leq 79)
\]
\[
P(X \geq 80) = 1 - 0.8166 = 0.1834
\]
\newline
Final Answer:
\newline
The approximate probability that the 2500 customers will "tie up" the 80 trunk lines at any given time is 0.1834, or 18.34\%.
\newpage

\section*{Projects for Chapter 3}
\subsection*{3A Mixture distribution}

In statistical modeling, if the data are contaminated by outliers or if the samples are drawn from a population formed by a mixture of two populations, one could use mixture distributions. Mixture distributions are used frequently in medical applications, such as microarray analysis. Suppose a random variable \(X\) has pdf \(f_1(x)\) with probability \(p_1\) and pdf \(f_2(x)\) with probability \(p_2\), where \(p_1 + p_2 = 1\). Then we say that the random variable \(X\) has a \textit{mixture distribution}. This can be thought of as observing a Bernoulli random variable \(Z\) that is equal to 1 with probability \(p_1\) and 2 with probability \(p_2\). Thus,

\[
X =
\begin{cases}
X_1 \sim f_1(x), & \text{if } Y = 1, \\
X_2 \sim f_2(x), & \text{if } Y = 2.
\end{cases}
\]

\begin{enumerate}
    \item[(a)] Show that the pdf of \(X\) is given by \(f(x) = p_1f_1(x) + p_2f_2(x)\).
    
    \item[(b)] If \((\mu_1, \sigma_1^2)\) and \((\mu_2, \sigma_2^2)\) are means and variances of \(f_1(x)\) and \(f_2(x)\), respectively, show that
    \[
    \mu = E(X) = p_1\mu_1 + p_2\mu_2,
    \]
    and
    \[
    \sigma^2 = \text{Var}(X) = p_1\sigma_1^2 + p_2\sigma_2^2 + p_1\mu_1^2 + p_2\mu_2^2 - (p_1\mu_1 + p_2\mu_2)^2.
    \]
\end{enumerate}

\section*{Solution}

Solution

Part (a): Show that the pdf of \(X\) is given by \(f(x) = p_1f_1(x) + p_2f_2(x)\)

Since \(X\) is drawn from one of two distributions, depending on the outcome of a Bernoulli random variable \(Z\), the overall probability density function (pdf) of \(X\) can be described as a mixture of the two individual pdfs, weighted by their respective probabilities.

\begin{itemize}
    \item If \(Z = 1\), \(X \sim f_1(x)\) with probability \(p_1\).
    \item If \(Z = 2\), \(X \sim f_2(x)\) with probability \(p_2\).
\end{itemize}

The total pdf \(f(x)\) is the weighted sum of the two pdfs:
\[
f(x) = P(Z = 1)f_1(x) + P(Z = 2)f_2(x)
\]
Substituting \(P(Z = 1) = p_1\) and \(P(Z = 2) = p_2\):
\[
f(x) = p_1 f_1(x) + p_2 f_2(x)
\]

Thus, the pdf of \(X\) is \(f(x) = p_1f_1(x) + p_2f_2(x)\).

---

Part (b): Show that \( \mu = E(X) = p_1\mu_1 + p_2\mu_2 \) and \( \sigma^2 = \text{Var}(X) = p_1\sigma_1^2 + p_2\sigma_2^2 + p_1\mu_1^2 + p_2\mu_2^2 - (p_1\mu_1 + p_2\mu_2)^2 \)
\newline

\textbf{Step 1: Mean of \(X\)}

The expected value of \(X\) is the weighted sum of the expected values of \(f_1(x)\) and \(f_2(x)\):

\[
\mu = E(X) = E(X | Z = 1)P(Z = 1) + E(X | Z = 2)P(Z = 2)
\]
Substituting the expected values and probabilities:
\[
\mu = \mu_1 p_1 + \mu_2 p_2
\]

Thus, \( \mu = p_1\mu_1 + p_2\mu_2 \).

---

\textbf{Step 2: Variance of \(X\)}

The variance of \(X\) is given by:
\[
\text{Var}(X) = E[X^2] - (E[X])^2
\]

We calculate \(E[X^2]\) first:
\[
E[X^2] = E[X^2 | Z = 1]P(Z = 1) + E[X^2 | Z = 2]P(Z = 2)
\]
Now, \(E[X^2 | Z = 1] = \sigma_1^2 + \mu_1^2\) and \(E[X^2 | Z = 2] = \sigma_2^2 + \mu_2^2\):
\[
E[X^2] = (\sigma_1^2 + \mu_1^2)p_1 + (\sigma_2^2 + \mu_2^2)p_2
\]

Now calculate \(\text{Var}(X)\):
\[
\text{Var}(X) = E[X^2] - (E[X])^2
\]
Substituting \(E[X^2]\) and \(E[X]\):
\[
\text{Var}(X) = p_1(\sigma_1^2 + \mu_1^2) + p_2(\sigma_2^2 + \mu_2^2) - (p_1\mu_1 + p_2\mu_2)^2
\]
Expanding the terms:
\[
\text{Var}(X) = p_1\sigma_1^2 + p_2\sigma_2^2 + p_1\mu_1^2 + p_2\mu_2^2 - (p_1\mu_1 + p_2\mu_2)^2
\]

Thus:
\[
\sigma^2 = \text{Var}(X) = p_1\sigma_1^2 + p_2\sigma_2^2 + p_1\mu_1^2 + p_2\mu_2^2 - (p_1\mu_1 + p_2\mu_2)^2
\]

\subsection*{3F Generating a large amount of random data using R}

Although Projects 3B and 3E can be used for random sample generation, it is tedious for large amounts of data generation. Now we will give some R commands for generating \(n\) random samples from a few common distributions.\newline

\texttt{runif(n, min=a, max=b)} will generate \(n\) random values from \(U(a, b)\). For instance, to generate 100 random samples from uniform \(U(0, 2)\), use the command \texttt{runif(100, 0, 2)}.\newline

Similarly, \texttt{rnorm(n, mean=a, sd=b)} will generate \(n\) random values from normal distribution with mean \(a\) and standard deviation \(b\). Thus, \texttt{rnorm(100, mean=4, sd=2)} will generate 100 random values from a normal distribution with mean 4 and standard deviation 2. Similarly, \texttt{rbinom(n, m, p)} will generate \(n\) random values from a binomial distribution with parameters \(m\) and \(p\). For instance, to generate 100 sample values from binomial with \(m=10\), and \(p=0.5\), use \texttt{rbinom(100, 10, 0.5)}. Similarly, \texttt{rexp(n, rate)} will generate \(n\) random samples from an exponential distribution with a specified rate. Thus, \texttt{rexp(100, 1/1000)} will generate 100 exponential values with rate \(\lambda = 1/1000\). \newline

Similarly, we can generate random samples from any distribution.\newline


\textbf{Exercise 3F.} For each of the generated random samples, write an R-code to create a histogram overlapped by the corresponding density function.

\section*{Solution}

% Customization for code highlighting
\lstset{
    basicstyle=\ttfamily\small,
    keywordstyle=\color{blue}\ttfamily,
    stringstyle=\color{red}\ttfamily,
    commentstyle=\color{green}\ttfamily,
    backgroundcolor=\color{gray!10},
    showstringspaces=false,
    frame=single,
    breaklines=true
}


\section*{Python Code for Generating Random Samples and Plotting}

\begin{lstlisting}[language=Python]
import numpy as np
import matplotlib.pyplot as plt
import seaborn as sns
from scipy.stats import uniform, norm, binom, expon

# Set the number of samples
n = 100

# Generate and plot for uniform distribution U(0, 2)
uniform_samples = np.random.uniform(0, 2, n)
sns.histplot(uniform_samples, kde=True, stat="density", bins=10, color="skyblue", label="Histogram")
x = np.linspace(0, 2, 1000)
plt.plot(x, uniform.pdf(x, loc=0, scale=2), 'r-', label="Uniform PDF")
plt.title("Uniform Distribution U(0, 2)")
plt.legend()
plt.show()

# Generate and plot for normal distribution N(4, 2^2)
normal_samples = np.random.normal(4, 2, n)
sns.histplot(normal_samples, kde=True, stat="density", bins=10, color="skyblue", label="Histogram")
x = np.linspace(-2, 10, 1000)
plt.plot(x, norm.pdf(x, loc=4, scale=2), 'r-', label="Normal PDF")
plt.title("Normal Distribution N(4, 2^2)")
plt.legend()
plt.show()

# Generate and plot for binomial distribution Binomial(10, 0.5)
binomial_samples = np.random.binomial(10, 0.5, n)
sns.histplot(binomial_samples, kde=False, stat="density", bins=10, color="skyblue", label="Histogram")
x = np.arange(0, 11)
plt.plot(x, binom.pmf(x, 10, 0.5), 'r-', marker="o", label="Binomial PMF")
plt.title("Binomial Distribution Binomial(10, 0.5)")
plt.legend()
plt.show()

# Generate and plot for exponential distribution Exp(1/1000)
exponential_samples = np.random.exponential(1000, n)
sns.histplot(exponential_samples, kde=True, stat="density", bins=10, color="skyblue", label="Histogram")
x = np.linspace(0, 5000, 1000)
plt.plot(x, expon.pdf(x, scale=1000), 'r-', label="Exponential PDF")
plt.title("Exponential Distribution Exp(1/1000)")
plt.legend()
plt.show()
\end{lstlisting}

\begin{figure}[H] % The H option forces the figure to be placed exactly here
\centering
% Replace this with the correct path to your Table 1.9 image file
\includegraphics[width=1\textwidth]{D:/Masters NPU Semeter/3rd Semester/Statistics/Homework/Submit/Homework 3/figs/Uniform Distribution.png}
\caption{Table  - Uniform Distribution.}
\end{figure}

\begin{figure}[H] % The H option forces the figure to be placed exactly here
\centering
% Replace this with the correct path to your Table 1.9 image file
\includegraphics[width=1\textwidth]{D:/Masters NPU Semeter/3rd Semester/Statistics/Homework/Submit/Homework 3/figs/Normal Distribution.png}
\caption{Table  - Normal Distribution.}
\end{figure}

\begin{figure}[H] % The H option forces the figure to be placed exactly here
\centering
% Replace this with the correct path to your Table 1.9 image file
\includegraphics[width=1\textwidth]{D:/Masters NPU Semeter/3rd Semester/Statistics/Homework/Submit/Homework 3/figs/Binomial Distribution.png}
\caption{Table  - Binomial Distribution.}
\end{figure}

\begin{figure}[H] % The H option forces the figure to be placed exactly here
\centering
% Replace this with the correct path to your Table 1.9 image file
\includegraphics[width=1\textwidth]{D:/Masters NPU Semeter/3rd Semester/Statistics/Homework/Submit/Homework 3/figs/Exponential Distribution.png}
\caption{Table  - Exponential Distribution.}
\end{figure}

\end{document}